% Created 2026-10-03 sáb 10:30
% Intended LaTeX compiler: xelatex
\documentclass[11pt]{article}

\usepackage{amsmath}
\usepackage{fontspec}
\usepackage{graphicx}
\usepackage{longtable}
\usepackage{wrapfig}
\usepackage{rotating}
\usepackage[normalem]{ulem}
\usepackage{capt-of}
\usepackage{hyperref}
\author{Nícolas Rosenthal Dal Corso}
\date{2026}
\title{Implementação do Simply Typed Lambda Calculus em Rust}
\hypersetup{
 pdfauthor={Nícolas Rosenthal Dal Corso},
 pdftitle={Implementação do Simply Typed Lambda Calculus em Rust},
 pdfkeywords={},
 pdfsubject={},
 pdfcreator={},
 pdflang={English}}
\begin{document}

\maketitle
\tableofcontents

\section{Introdução}
\label{sec:org412601e}

Este relatório apresenta a implementação, em Rust, do \textbf{Simply Typed Lambda}
\textbf{Calculus} (STLC), seguindo a formalização apresentada por Benjamin C. Pierce
em \textbf{Types and Programming Languages} (TAPL).

O cálculo lambda simplesmente tipado constitui uma extensão do cálculo lambda
não tipado na qual cada parâmetro de uma função recebe um tipo explícito.
Além das abstrações e aplicações características do cálculo lambda, a
linguagem considerada nesta implementação possui valores booleanos e
expressões condicionais. Dessa forma, a implementação contempla os elementos
fundamentais necessários para estudar a relação entre sintaxe, tipagem e
avaliação de uma linguagem funcional tipada.

A implementação foi desenvolvida em Rust, utilizando tipos algébricos
representados por \emph{enums} e estruturas recursivas por meio de
\texttt{Box<T>}. Essa representação permite modelar diretamente a estrutura
abstrata dos termos da linguagem e, ao mesmo tempo, aproveitar o sistema de
tipos da própria linguagem Rust para garantir propriedades estruturais da
implementação.

A linguagem é organizada em torno de três componentes principais: os
\textbf{termos}, que definem a sintaxe das expressões; os \textbf{tipos}, que descrevem as
categorias sintáticas às quais os termos pertencem; e as regras de avaliação
e tipagem, que determinam quais expressões são válidas e como elas são
reduzidas.

Neste relatório, a implementação é apresentada a partir de sua
correspondência com a definição formal da linguagem. Inicialmente são
descritos os termos e sua representação em Rust; em seguida, podem ser
apresentados os tipos, as regras de tipagem e a avaliação operacional da
linguagem.
\section{Termos}
\label{sec:org1cbfd93}

Os termos da linguagem correspondem às expressões que podem ser construídas
pelo STLC. A implementação considera seis formas sintáticas fundamentais:
\section{variáveis;}
\label{sec:orgd5eff92}
\section{o valor booleano \texttt{true};}
\label{sec:org8abbb0c}
\section{o valor booleano \texttt{false};}
\label{sec:orgb34ae02}
\section{expressões condicionais;}
\label{sec:org0f1daee}
\section{abstrações lambda;}
\label{sec:orgdd7b30b}
\section{aplicações de funções.}
\label{sec:orgc6ebb9b}

Em Rust, essas formas são representadas pelo tipo algébrico \texttt{Term}:

\begin{verbatim}
use std::collections::BTreeSet;
use std::fmt;

use super::types::Type;

#[derive(Debug, Clone, PartialEq, Eq)]
pub enum Term {
/// `x`
Var(String),
/// `true`
True,
/// `false`
False,
/// `if c then t else e`
If { condition: Box<Term>, then_branch: Box<Term>, else_branch: Box<Term> },
/// `λx:T. body`
Lambda { param: String, ty: Type, body: Box<Term> },
/// `f x`
App { func: Box<Term>, arg: Box<Term> },
}
\end{verbatim}

A enumeração \texttt{Term} constitui uma representação direta da sintaxe abstrata
da linguagem. Cada variante corresponde a uma construção sintática
específica.
\subsection{Sintaxe concreta}
\label{sec:orgcc7fcbd}

A sintaxe dos termos pode ser descrita pela seguinte gramática BNF:

{[}
\begin{array}{rcl}
t &::=& x \
&\mid& \mathrm{true} \
&\mid& \mathrm{false} \
&\mid& \mathrm{if}\ t\ \mathrm{then}\ t\ \mathrm{else}\ t \
&\mid& \lambda x.\ t \
&\mid& t\ t
\end{array}
]

Em notação BNF, a mesma gramática pode ser apresentada como:

{[}
\begin{array}{rcl}
t &::=& x
\mid \mathrm{true}
\mid \mathrm{false}
\mid \mathrm{if}\ t\ \mathrm{then}\ t\ \mathrm{else}\ t
\mid \lambda x.\ t
\mid t\ t
\end{array}
]

onde \(t\) representa um termo, \(x\) representa uma variável e \(T\) representa
um tipo.

A gramática estabelece apenas quais expressões pertencem à sintaxe da
linguagem. A validade tipada dessas expressões é determinada posteriormente
pelas regras de tipagem.
\subsection{Variáveis}
\label{sec:orgb66cdd6}

Uma variável é representada pela produção

{[}
t ::= x.
]

Na implementação, essa construção corresponde à variante:

\begin{verbatim}
Var(String)
\end{verbatim}

O identificador da variável é armazenado como uma \texttt{String}. Assim, por
exemplo, a expressão matemática

{[}
x
]

pode ser representada por:

\begin{verbatim}
Term::Var("x".to_string())
\end{verbatim}

A utilização de \texttt{String}, em vez de índices inteiros ou de uma estrutura
separada para identificadores, mantém a representação próxima da notação
utilizada na apresentação formal do cálculo.
\subsection{Valores booleanos}
\label{sec:org12aacda}

A linguagem possui dois valores booleanos:

{[}
\mathrm{true}
\qquad
\mathrm{false}.
]

Eles são representados diretamente pelas variantes:

\begin{verbatim}
True,
False,
\end{verbatim}

correspondendo respectivamente a:

{[}
\mathrm{true} : \mathrm{Bool}
]

e

{[}
\mathrm{false} : \mathrm{Bool}.
]
\subsection{Expressões condicionais}
\label{sec:org3908273}

Uma expressão condicional possui a forma

{[}
\mathrm{if}$\backslash$ t\_1$\backslash$ \mathrm{then}$\backslash$ t\_2$\backslash$ \mathrm{else}$\backslash$ t\_3.
]

Na implementação, ela é representada por:

\begin{verbatim}
If {
condition: Box<Term>,
then_branch: Box<Term>,
else_branch: Box<Term>,
}
\end{verbatim}

Os três campos são termos porque tanto a condição quanto os dois ramos
podem conter expressões arbitrariamente complexas.

Por exemplo,

{[}
\mathrm{if}$\backslash$ \mathrm{true}$\backslash$ \mathrm{then}$\backslash$ x$\backslash$ \mathrm{else}$\backslash$ y
]

pode ser representado por:

\begin{verbatim}
Term::If {
condition: Box::new(Term::True),
then_branch: Box::new(Term::Var("x".to_string())),
else_branch: Box::new(Term::Var("y".to_string())),
}
\end{verbatim}

A utilização de \texttt{Box<Term>} é necessária porque \texttt{Term} é um tipo recursivo:
um termo condicional contém outros termos.
\subsection{Abstrações lambda}
\label{sec:org37a4b88}

A abstração lambda constitui a construção utilizada para representar
funções. No STLC, diferentemente do cálculo lambda não tipado, o parâmetro
da função possui um tipo explícito:

{[}
\(\lambda\) x.$\backslash$ t.
]

A representação em Rust é:

\begin{verbatim}
Lambda {
param: String,
ty: Type,
body: Box<Term>,
}
\end{verbatim}

O campo \texttt{param} armazena o nome do parâmetro, enquanto \texttt{ty} armazena seu
tipo. O corpo da função é novamente um \texttt{Term} e, portanto, é armazenado
indiretamente através de \texttt{Box<Term>}.

Por exemplo, a função identidade sobre booleanos,

{[}
\(\lambda\) x:\mathrm{Bool}.x,
]

pode ser representada por:

\begin{verbatim}
Term::Lambda {
param: "x".to_string(),
ty: Type::Bool,
body: Box::new(Term::Var("x".to_string())),
}
\end{verbatim}

A presença do tipo no próprio termo permite representar diretamente a
sintaxe anotada do STLC.
\subsection{Aplicações}
\label{sec:orgeb4ea92}

A aplicação de uma função a um argumento possui a forma

{[}
t\_1$\backslash$ t\_2.
]

Ela é representada pela variante:

\begin{verbatim}
App {
func: Box<Term>,
arg: Box<Term>,
}
\end{verbatim}

Por exemplo, a aplicação da função identidade booleana ao valor \texttt{true},

{[}
(\(\lambda\) x:\mathrm{Bool}.x)$\backslash$ \mathrm{true},
]

pode ser representada por:

\begin{verbatim}
Term::App {
func: Box::new(Term::Lambda {
param: "x".to_string(),
ty: Type::Bool,
body: Box::new(Term::Var("x".to_string())),
}),
arg: Box::new(Term::True),
}
\end{verbatim}

Assim como ocorre com as demais construções recursivas, \texttt{Box<Term>} permite
que uma aplicação contenha termos arbitrariamente complexos tanto na posição
de função quanto na posição de argumento.
\subsection{Correspondência entre a sintaxe e a implementação}
\label{sec:org944c63a}

A correspondência entre a gramática e a representação em Rust pode ser
resumida da seguinte forma:

{[}
\begin{array}{c|c}
\text{Construção sintática} & \text{Representação Rust} \ \hline
x & \mathrm{Term::Var(String)} \
\mathrm{true} & \mathrm{Term::True} \
\mathrm{false} & \mathrm{Term::False} \
\mathrm{if}\ t_1\ \mathrm{then}\ t_2\ \mathrm{else}\ t_3
& \mathrm{Term::If} \
\lambda x.\ t
& \mathrm{Term::Lambda} \
t_1\ t_2
& \mathrm{Term::App}
\end{array}
]

Essa correspondência é intencionalmente próxima da definição formal. A
estrutura recursiva de \texttt{Term} reproduz a estrutura da árvore sintática
abstrata: cada termo composto contém subtermos representados por outros
valores de \texttt{Term}.

A derivação dos tipos desses termos, bem como as regras de avaliação, pode
ser definida independentemente da representação concreta utilizada pelo
programa. Isso permite separar a definição formal da linguagem de sua
implementação e utilizar a enumeração \texttt{Term} como uma representação
computacional da sintaxe abstrata.
\end{document}
